% !TEX root = ../../../main.tex
% 所属：第一部分：单变量分析学 / 第10章：一些定积分的问题
% 节：10.5 中值定理
% 正文按原文件顺序保存；由 ../chapter10.tex 调用。
\section{中值定理}

本节主要整理几道利用积分中值定理与 Rolle 中值定理的题目.

\begin{problem}
 设 $f(x)$ 在 $(0,1)$ 上可微，且满足
	\[
	f(1)-2\int_0^{\frac12}x f(x)\dd x=0.
	\]
	求证：在 $(0,1)$ 内至少有一点 $\xi$，使得
	\[
	f'(\xi)=-\frac{f(\xi)}{\xi}.
	\]
\end{problem}

\begin{proof}
	由积分中值定理，存在 $\eta\in\left(0,\dfrac12\right)$，使得
	\[
	\int_0^{\frac12}x f(x)\dd x=\frac12\eta f(\eta).
	\]
	因此
	\[
	f(1)=2\int_0^{\frac12}x f(x)\dd x=\eta f(\eta).
	\]
	令 $F(x)=x f(x)$，则
	\[
	F(1)=f(1)=\eta f(\eta)=F(\eta).
	\]
	于是对 $F(x)$ 在 $[\eta,1]$ 上应用 Rolle 中值定理，存在 $\xi\in(\eta,1)$，使得
	\[
	F'(\xi)=0.
	\]
	即
	\[
	f(\xi)+\xi f'(\xi)=0,
	\]
	所以
	\[
	f'(\xi)=-\frac{f(\xi)}{\xi}.
	\]
\end{proof}

\begin{problem}
设 $f(x)$ 在区间 $[0,1]$ 上连续，在 $(0,1)$ 上可导，满足
	\[
	f(1)=2\int_0^{\frac12}e^{1-x^2}f(x)\dd x.
	\]
	求证：存在一点 $\xi\in(0,1)$，使得
	\[
	f'(\xi)=2\xi f(\xi).
	\]
\end{problem}

\begin{proof}
	由题设两边同乘 $e^{-1}$ 得
	\[
	e^{-1}f(1)=2\int_0^{\frac12}e^{-x^2}f(x)\dd x.
	\]
	由积分中值定理，存在 $\eta\in\left(0,\dfrac12\right)$，使得
	\[
	2\int_0^{\frac12}e^{-x^2}f(x)\dd x=e^{-\eta^2}f(\eta).
	\]
	于是
	\[
	e^{-1}f(1)=e^{-\eta^2}f(\eta).
	\]
	令
	\[
	F(x)=e^{-x^2}f(x),
	\]
	则 $F(1)=F(\eta)$，由 Rolle 中值定理，存在 $\xi\in(\eta,1)$，使得
	\[
	F'(\xi)=0.
	\]
	又
	\[
	F'(x)=e^{-x^2}\left(f'(x)-2x f(x)\right),
	\]
	故
	\[
	f'(\xi)-2\xi f(\xi)=0,
	\]
	即
	\[
	f'(\xi)=2\xi f(\xi).
	\]
\end{proof}

\begin{problem}
设 $f(x)$ 在 $[0,1]$ 上连续，在 $(0,1)$ 内三次可微，并且满足条件：
	\[
	\int_0^{\frac14}f(t)\dd t=\frac14 f(0),
	\qquad
	\int_{\frac14}^{1}f(t)\dd t=\frac34 f(1).
	\]
	试证：若 $f(0)=f(1)$，则存在 $\xi\in(0,1)$，使得
	\[
	f'''(\xi)=0.
	\]
\end{problem}

\begin{proof}
	由积分中值定理，存在 $\xi_1\in\left(0,\dfrac14\right)$，使得
	\[
	\int_0^{\frac14}f(t)\dd t=\frac14 f(\xi_1).
	\]
	结合题设得
	\[
	f(\xi_1)=f(0).
	\]
	同理，存在 $\xi_2\in\left(\dfrac14,1\right)$，使得
	\[
	\int_{\frac14}^{1}f(t)\dd t=\frac34 f(\xi_2),
	\]
	从而
	\[
	f(\xi_2)=f(1).
	\]
	又 $f(0)=f(1)$，所以
	\[
	f(0)=f(\xi_1)=f(\xi_2)=f(1).
	\]
	于是分别在 $[0,\xi_1]$，$[\xi_1,\xi_2]$，$[\xi_2,1]$ 上应用 Rolle 中值定理，存在
	\[
	\eta_1\in(0,\xi_1),\qquad
	\eta_2\in(\xi_1,\xi_2),\qquad
	\eta_3\in(\xi_2,1),
	\]
	使得
	\[
	f'(\eta_1)=f'(\eta_2)=f'(\eta_3)=0.
	\]
	再对 $f'(x)$ 分别在 $[\eta_1,\eta_2]$ 与 $[\eta_2,\eta_3]$ 上应用 Rolle 中值定理，存在
	\[
	\delta_1\in(\eta_1,\eta_2),
	\qquad
	\delta_2\in(\eta_2,\eta_3),
	\]
	使得
	\[
	f''(\delta_1)=f''(\delta_2)=0.
	\]
	最后对 $f''(x)$ 在 $[\delta_1,\delta_2]$ 上应用 Rolle 中值定理，存在 $\xi\in(\delta_1,\delta_2)$，使得
	\[
	f'''(\xi)=0.
	\]
\end{proof}

\begin{problem}
	设函数 $f(x)$ 在 $[0,\pi]$ 上连续，且
	\[
	\int_0^\pi f(x)\dd x=0,
	\qquad
	\int_0^\pi f(x)\cos x\dd x=0.
	\]
	证明：在 $(0,\pi)$ 内至少存在两个不同的点 $\xi_1,\xi_2$，使得
	\[
	f(\xi_1)=f(\xi_2)=0.
	\]
\end{problem}

\begin{proof}
	令
	\[
	g(x)=\int_0^x f(t)\dd t,
	\qquad
	h(x)=\int_0^x g(t)\sin t\dd t.
	\]
	由题设知
	\[
	g(0)=g(\pi)=0,
	\qquad
	h(0)=0.
	\]
	又由分部积分，
	\[
	\begin{aligned}
		h(\pi)
		&=\int_0^\pi g(x)\sin x\dd x \\
		&=\left[-g(x)\cos x\right]_0^\pi+
		\int_0^\pi f(x)\cos x\dd x=0.
	\end{aligned}
	\]
	于是 $h(0)=h(\pi)=0$，由 Rolle 中值定理，存在 $\eta\in(0,\pi)$，使得
	\[
	h'(\eta)=g(\eta)\sin\eta=0.
	\]
	由于 $\eta\in(0,\pi)$，故 $\sin\eta>0$，从而
	\[
	g(\eta)=0.
	\]
	因此
	\[
	g(0)=g(\eta)=g(\pi)=0.
	\]
	分别对 $g(x)$ 在 $[0,\eta]$ 与 $[\eta,\pi]$ 上应用 Rolle 中值定理，存在
	\[
	\xi_1\in(0,\eta),
	\qquad
	\xi_2\in(\eta,\pi),
	\]
	使得
	\[
	g'(\xi_1)=g'(\xi_2)=0.
	\]
	又 $g'(x)=f(x)$，所以
	\[
	f(\xi_1)=f(\xi_2)=0.
	\]
	显然 $\xi_1\neq\xi_2$.
\end{proof}

